% SEGMENT OLP-0355-S01
% Part: lambda-calculus
% Chapter: introduction
% Section: minimization

\documentclass[../../../include/open-logic-section]{subfiles}

\begin{document}

\olfileid{lam}{int}{min}
\olsection{\usetoken{S}{lambda definable} சார்புகள் சிறுமமாக்கலின் கீழ் மூடியவை}

\begin{lem}
$f(x,y)$ ஆனது !!{lambda definable} என்று கொள்க. $g$ ஐ
\[
g(x) \simeq \umin{y}{f(x,y)}.
\]
என வரையறுக்க. அப்போது $g$ ஆனது !!{lambda definable}.
\end{lem}

\begin{proof}
தோராயமான கருத்து பின்வருமாறு. $x$ கொடுக்கப்பட்டால்,~$n$ இலிருந்து தொடங்கி
ஒரு~$y$ ஐத் தேடும் சார்பு $h_x(n)$ ஐ வரையறுக்க நிலைப்புள்ளி லாம்டா
சொல்~$Y$ ஐப் பயன்படுத்துவோம்; பின்னர் $g(x)$ என்பது $h_x(0)$ மட்டுமே.
சார்பு~$h_x$ ஐ ஒரு நிலைப்புள்ளிச் சமன்பாட்டின் தீர்வாக எழுதலாம்:
\[
h_x(n) \simeq
\begin{cases}
n & \text{$f(x,n) = 0$ எனில்} \\
h_x(n+1) & \text{இல்லையெனில்.}
\end{cases}
\]

விவரங்கள் பின்வருமாறு. $f$ லாம்டா-வரையறுக்கத்தக்கது என்ற கருதுகோளால்,
% SOURCE CORRECTION TA-LAM-010: the lemma assumes lambda-definability, not primitive recursiveness.
அது ஏதோ ஒரு சொல் $F$ ஆல் !!{lambda defined} ஆக உள்ளது. $D(M, N, \bar{0})
\red M$ மற்றும் $D(M, N, \bar{1}) \red N$ ஆக அமையும் ஒரு லாம்டா சொல்~$D$
நம்மிடம் இருப்பதையும் நினைவுகொள்க. இப்போதைக்கு $x$ ஐ நிலையாகக் கொண்டு,
$h_x$ ஐ !!{lambda define} செய்ய, பின்வருவதை நிறைவுசெய்யும் ஒரு சொல்~$H$
(அது~$x$ ஐச் சார்ந்தது) நமக்குத் தேவை:
\[
H(\num{n}) \equiv D(\num{n}, H(S(\num{n})), F(x, \num{n})).
\]
நிலைப்புள்ளிச் சொல்~$Y$ ஐப் பயன்படுத்தி இதைச் செய்யலாம். முதலில் $U$ என்பது
\[
\lambd[h][\lambd[z][D(z,(h(Sz)),F(x,z))]],
\]
என்ற சொல்லாக இருக்கட்டும்; பின்னர் $H$ என்பது~$YU$ என்ற சொல்லாக இருக்கட்டும்.
~$H$ இல் உள்ள ஒரே கட்டற்ற மாறி~$x$ என்பதைக் கவனிக்க. மேலுள்ள சமன்பாட்டை
$H$ நிறைவுசெய்கிறது என்று காட்டுவோம்.

$Y$ இன் வரையறையால்,
\[
H = YU \equiv U(YU) = U(H).
\]
குறிப்பாக, ஒவ்வோர் இயல் எண்~$n$ க்கும்
\begin{align*}
H(\num{n}) & \equiv U(H, \num{n}) \\
& \red D(\num{n}, H(S(\num{n})), F(x, \num{n})),
\end{align*}
தேவையானவாறே கிடைக்கிறது. கடைசி வரியில்~$x$ க்குப் பதிலாக ஓர் எண்ணுரு
$\num{m}$ ஐ இடைநிறுத்தினால், $F(\num{m}, \num{n})$ ஆனது $\num{0}$ குக்
குறைந்தால் கோவை $\num{n}$ குக் குறைகிறது; $F(\num{m}, \num{n})$ வேறு
ஏதேனும் எண்ணுருவிற்குக்
குறைந்தால் கோவை $H(S(\num{n}))$ குக் குறைகிறது என்பதைக் கவனிக்க.

நிறுவலை முடிக்க, $G$ என்பது $\lambd[x][H(\num 0)]$ ஆக இருக்கட்டும். அப்போது
$G$ ஆனது~$g$ ஐ !!{lambda define}s; வேறு சொற்களில், ஒவ்வொரு~$m$ க்கும்
$g(m)$ வரையறுக்கப்பட்டால் $G(\num m)$ ஆனது~$\overline {g(m)}$ குக் குறைகிறது;
இல்லையெனில் அதற்கு இயல்புவடிவம் இல்லை.
\end{proof}

\end{document}
